How to determine the set of arbitrage-free prices in the discrete time setting
for a continuously distributed underlying?

The support of a Borel probability measure \( \nu \) on \( \mathbb{R}^d \) is the smallest
closed set \( A \subset \mathbb{R}^d \) such that \( \nu(A^c) = 0 \), and it will be denoted by \( \text{supp } \nu \).
It holds
\( \text{supp } \nu = \bigcap_{\nu(A^c)=0, A \text{ closed}} A \).

Further let the convex hull of the support be
\( \Gamma\ \mu := \text{conv } \text{supp } \mu = \{ \sum_{k=1}^n \alpha_k y_k | \alpha_k \geq 0 , \sum_k \alpha_k = 1, y_k \in \text{ supp } \mu \} \)

This is usually determined by finding the extreme points (those points which don't are inside (but can be an endpoint) a line connecting two points in \( C \)) and connecting them in a manner such that no intersections occur.

The relative interior is defined for a convex set \( C \subset \mathbb{R}^d \) as the set of \( x \in C \), s.t. for all \( y \in C \), 
there is some \( \varepsilon > 0 \) with \( x - \epsilon(y - x) \in C \). The notation is \( \text{ri } C \).

With considerable effort one can then derive: 
Theorem:
Let \( \mu \) be the distribution of the discounted price vector \( \frac{S}{(1 + r)} \) of
the risky assets. Then the market model is arbitrage-free if and only if the price system \( \pi \) belongs 
to the relative interior \( \text{ri }\Gamma\ \mu \) of the convex hull of the support of \( \mu \).

Suppose \( (S^0, S^1, S^2) = \pi \) and \( S^2 \) does not satisfy the conditions above, then an arbitrage strategy is usually given
by going long/short on \( S^2 \) (depending on if it's lower, higher than the arbitrage free price respectively) and doing the opposite 
with the proceeds to \( S^1 \) s.t. the resulting portfolio is at least the upper envelope of \( \text{ri }\Gamma\ \mu \).

If this is supposed to be minimal it needs to be exactly the upper envelope for all points. If \( S^2 \) is not linearly dependent on \( S^1 \)
(that is we can't use a simple trading strategy in terms of \( S^1, S^0 \) to replicate \( S^2 \)), then we need to use the maximal linear part 
of the upper envelope.


Derive the price of a Chooser Option

At \( T_1 \) the owner may choose to obtain an european Call or Put both with strike \( K \) and maturity \( T_2 \).
Under the assumption that the holder of the option always chooses the option with higher value at \( T_1 \), we get
via FTAP
\( V_0 = \frac{1}{(1 + r)^T} \mathbb{E}[(S^1_{T_2} - K)_{+}|\mathcal{F}_0] + \mathbb{E}[K - S_{T_2}|\mathcal{F}_0, (C_{T_2,K})_{T_1} < (P_{T_2,K})_{T_1}] \)
and with the Put-Call parity \( (C_{T_2,K})_{T_1} < (P_{T_2,K})_{T_1} \iff 0 < (P_{T_2,K})_{T_1} - (C_{T_2,K})_{T_1} = (1 + r)^{T_1 - T_2}K - S_{T_1} \)
inserting it with \( A := \{ 0 < (1 + r)^{T_1 - T_2}K - S_{T_1} \} \) yields
\( V_0 = (C_{T_2,K})_{0} + \frac{1}{(1+r)^T}\mathbb{E}[K - S_{T_2}|\mathcal{F}_0, A] \)
Using on the latter summand the Tower-property 
\( \frac{1}{(1+r)^T}\mathbb{E}[K - S_{T_2}|\mathcal{F}_0, A] = \mathbb{E}[\frac{1}{(1+r)^T}\mathbb{E}[K - S_{T_2}|\mathcal{F}_{T_1}]|\mathcal{F}_0, A] \)
and applying FTAP on \( S_{T_2} \) simplifies this to
\( \mathbb{E}[\frac{1}{(1+r)^{T_2}}K - \frac{1}{(1+r)^{T_2}}S_{T_1}|\mathcal{F}_0, A] \)
\( = \frac{1}{(1+r)^{T_1}} \mathbb{E}[(1 + r)^{T_1 - T_2}K - S_{T_1}|\mathcal{F}_0, A] \)
applying what we know from the conditioning on \( A \)
\( = \frac{1}{(1+r)^{T_1}} \mathbb{E}[((1 + r)^{T_1 - T_2}K - S_{T_1})_{+}|\mathcal{F}_0] = (P_{T_1, (1+r)^{T_1 - T_2}K})_0 \).
Summarizing 

\( V_0 = (C_{T_2,K})_{0} + (P_{T_1, (1+r)^{T_1 - T_2}K})_0 \)

Prove an equivalent characterization for self-financing trading strategies:

To demonstrate that a trading strategy \(\vartheta\) is self-financing under the condition \(\sum_{i=0}^d S^i_k(\vartheta^i_k - \vartheta^i_{k-1}) = 0\), we need to show that the change in value of the portfolio solely depends on changes in asset prices, rather than any injection or withdrawal of funds.
The value of the portfolio at time \(t\) \(V_t\), is given by:

\[
V_t = \sum_{i=0}^n \vartheta^i_t S^i_t.
\]

A trading strategy is self-financing if the change in the portfolio value from time \(k-1\) to \(k\) is due only to changes in asset prices, i.e.,

\[
V_k - V_{k-1} = \sum_{i=0}^d \vartheta^i_{k-1} (S^i_k - S^i_{k-1}).
\]

Let's analyze the change in value of the portfolio:

\[
V_k - V_{k-1} = \sum_{i=0}^n \vartheta^i_k S^i_k - \sum_{i=0}^n \vartheta^i_{k-1} S^i_{k-1}.
\]

We can rewrite this as:

\[
V_k - V_{k-1} = \sum_{i=0}^n (\vartheta^i_k S^i_k - \vartheta^i_{k-1} S^i_{k-1}).
\]

This can be decomposed into:

\[
V_k - V_{k-1} = \sum_{i=0}^n \vartheta^i_k S^i_k - \sum_{i=0}^n \vartheta^i_{k-1} S^i_k + \sum_{i=0}^n \vartheta^i_{k-1} S^i_k - \sum_{i=0}^n \vartheta^i_{k-1} S^i_{k-1}.
\]

Combining terms, we get:

\[
V_k - V_{k-1} = \sum_{i=0}^n (\vartheta^i_k - \vartheta^i_{k-1}) S^i_k + \sum_{i=0}^n \vartheta^i_{k-1} (S^i_k - S^i_{k-1}).
\]

Given the condition \(\sum_{i=0}^d S^i_k (\vartheta^i_k - \vartheta^i_{k-1}) = 0\), the first term vanishes. Thus we are left with:

\[
V_k - V_{k-1} = \sum_{i=0}^n \vartheta^i_{k-1} (S^i_k - S^i_{k-1}).
\]

This shows that the change in the portfolio's value is entirely due to changes in asset prices, rather than changes in portfolio positions. Therefore, the strategy \(\vartheta\) is self-financing.


Explain capital allocation principles and state examples

Mathematical description
Consider \(n\) risky sub-portfolios whose outcomes are given by \(X_1, \ldots, X_n\).
Then the total portfolio is given by \(X := X_1 + \cdots + X_n\).
If risk is measured in terms of a sub-additive risk measure \(\rho\), then there is a diversification effect in the total portfolio:
\(
\rho(X) \leq \sum_{i=1}^n \rho(X_i)
\)

Goal
Allocate to each sub-portfolio \(X_i\), \(i = 1, \ldots, n\), a risk capital contribution \(\rho(X_i \mid X)\) which reflects its contribution 
to the total risk \(\rho(X)\) and to the diversification effect.

Key properties of a meaningful capital allocation
Complete capital allocation: \(\sum_{i=1}^n \rho(X_i \mid X) = \rho(X)\)
Diversification: For each sub-portfolio \(X_i\), \(i = 1, \ldots, n\), the allocated capital is at most as large as the sub-portfolio’s standalone capital:
\(
\rho(X_i \mid X) \leq \rho(X_i) \quad \text{for } i = 1, \ldots, n.
\)


Proportional allocation:
\(
\rho(X_i \mid X) = \frac{\rho(X_i)}{\sum_{j=1}^n \rho(X_j)} \rho(X)
\)

Marginal principles, e.g.:
Euler allocation (for homogeneous risk measures):
\(
\rho(X_i \mid X) = \left. \frac{d}{dh} \rho(X + h X_i) \right|_{h=0}
\)
 Covariance principle:
\(
\rho(X_i \mid X) = \frac{\text{Cov}(X_i, X)}{\sigma^2(X)} \rho(X)
\)

Euler's capital allocation principle can be expressed more succinctly if \( \rho(X) = f(X) \)
for some \( f : \mathbb{R}^n \to \mathbb{R} \) which is positive homogeneous of order \( \alpha \).
It holds 
\( \rho(X_i|X) = \frac{\partial f_{\rho}}{\partial x_i}(\mathbf{1}) \)
and it satisfies both Complete capital allocation and Diversification.
Under certain technical conditions the \( \text{V@R}_{\lambda} \) actually satisfies this with 
\( \text{V@R}_{\lambda} = \mathbb{E}[-X_i| -X =\text{V@R}_{\lambda}(X)] \)
